% TOP_PROBLEM: {"id": 10900041, "title": "5.16 (Walsh) — Hyperbolic groups with Kleinian-type boundaries", "category_id": 17, "category": {"id": 17, "name": "group_theory", "display_name": "Group Theory", "description": "Problems about groups, group actions, representations, and related algebraic structures.", "slug": "group-theory", "order_index": 17, "created_at": "2026-07-31T15:26:25.671Z"}, "status": "partially_solved"}
% FIRST_POSTED: 2026-10-06
% ATTEMPT_STATUS: unresolved
% OPENAI_CLAIM: {"schema_version":1,"source_repo":"https://github.com/openai/math","source_commit":"adc7f1241b42e322a6451854ab7e4b4c146bf78a","release_date":"2026-10-06","match":"partial","resolution":"partial","independently_reviewed":false,"second_pass":"n/a","families":[{"family":"246","title":"Cannon's conjecture","subject":"Group theory","manuscripts":[{"title":"A Modulus Proof of Cannon’s Conjecture","pdf_path":"preprints/A-Modulus-Proof-of-Cannons-Conjecture-September-23-2026/paper.pdf","readme_path":"preprints/A-Modulus-Proof-of-Cannons-Conjecture-September-23-2026/README.md","date":"2026-09-23","theorem_ref":"Theorem 1.1"}],"lean":{"doc_path":"lean/docs/246.md","comparators":["lean/ComparatorChallenges/CannonGeometricAction.json","lean/ComparatorChallenges/CannonGeometricAction.lean"],"declarations":[],"covers_main_theorem":true},"reasoning_trace":null}]}
\documentclass[11pt]{article}
\usepackage[margin=25mm]{geometry}
\usepackage{fontspec}
\setmainfont{Latin Modern Roman}
\usepackage{amsmath,amssymb,mathtools}
\usepackage{unicode-math}
\setmathfont{Latin Modern Math}
\usepackage{enumitem}
\usepackage{xurl,hyperref,bookmark}
\hypersetup{colorlinks=true,urlcolor=blue}
\setcounter{secnumdepth}{0}
\setlength{\parindent}{0pt}
\setlength{\parskip}{5pt}
\setlength{\emergencystretch}{3em}
\begin{document}
\section*{5.16 (Walsh) — Hyperbolic groups with Kleinian-type boundaries}

\subsection{OpenAI's claim}
OpenAI presents a mathematical result relevant to this question in the selected manuscript. A hyperbolic group with two-sphere boundary is claimed to act geometrically on hyperbolic three-space with finite kernel.

\subsection{Outline of the argument}
The argument applies a boundary modulus criterion. Unbounded modulus produces scalar limits analyzed by exponential records or subexponential probability flows; comparing same-depth expansions forces the required measure control and the three-dimensional hyperbolic action.

\subsection{Scope relative to this problem}
The main theorem and its closed-hyperbolic-three-manifold consequences settle the S\textasciicircum{}2-boundary subcase. The record quantifies over additional boundary types, general hyperbolic groups, dimensions, or independent bundle clauses that this paper does not settle. This entry records partial support; complete coverage of the question is not established by this catalogue.

\subsection{Formal verification}
For family 246, the documented theorem gives a proper cocompact hyperbolic-three-space action with finite kernel from the two-sphere boundary hypothesis. The listed Comparator files are \texttt{CannonGeometricAction.lean}. The inventory does not supply a named declaration for these Comparator entries. This site has not executed the Lean code. The scope descriptions do not certify the remaining manuscript arguments.

\subsection{Status caveat}
OpenAI attributes these manuscripts to work by an internal model. Its release cautions that arguments without formalization can contain mistakes. This catalogue has not independently checked the proofs; the labels report the mathematical claims attributed to OpenAI.

\subsection{Sources}
\begin{itemize}
\item OpenAI, \emph{A Modulus Proof of Cannon’s Conjecture}, 2026-09-23; citation key \texttt{OAI:A-Modulus-Proof-of-Cannons-Conjecture-September-23-2026}. \href{https://github.com/openai/math/blob/adc7f1241b42e322a6451854ab7e4b4c146bf78a/preprints/A-Modulus-Proof-of-Cannons-Conjecture-September-23-2026/paper.pdf}{PDF (GitHub)}; \href{https://github.com/openai/math/raw/adc7f1241b42e322a6451854ab7e4b4c146bf78a/preprints/A-Modulus-Proof-of-Cannons-Conjecture-September-23-2026/paper.pdf}{PDF (direct)}; \href{https://github.com/openai/math/blob/adc7f1241b42e322a6451854ab7e4b4c146bf78a/preprints/A-Modulus-Proof-of-Cannons-Conjecture-September-23-2026/README.md}{Manuscript page and citation}; \href{https://github.com/openai/math/blob/main/preprints/A-Modulus-Proof-of-Cannons-Conjecture-September-23-2026/paper.pdf}{Latest version}.
\item Family 246: \href{https://github.com/openai/math/blob/adc7f1241b42e322a6451854ab7e4b4c146bf78a/lean/docs/246.md}{Lean scope}.
\item \href{https://github.com/openai/math/blob/adc7f1241b42e322a6451854ab7e4b4c146bf78a/lean/ComparatorChallenges/CannonGeometricAction.json}{Comparator CannonGeometricAction.json}.
\item \href{https://github.com/openai/math/blob/adc7f1241b42e322a6451854ab7e4b4c146bf78a/lean/ComparatorChallenges/CannonGeometricAction.lean}{Comparator CannonGeometricAction.lean}.
\item \href{https://github.com/openai/math/blob/adc7f1241b42e322a6451854ab7e4b4c146bf78a/CONTENTS.md}{OpenAI catalogue entry}.
\end{itemize}
\end{document}
